% © 2026 Ing. David Jaroš — CC BY-NC-ND 4.0
%
% File: research_tracks/T3_ALPHA/theta0_uv_sign_from_psi_casimir_candidate.tex
% Purpose: Candidate UV closure of the last alpha-proof condition: derive the
%          electroweak symmetry-breaking sign from the compact psi-sector.
% Status: [L2 candidate / conditional].  This is the strongest honest closure
%         without a fully specified canonical V(Theta).  It identifies the
%         exact extra UV hypotheses needed to promote the result to canonical.
% Date: 2026-06-13

% UBT-AI-PROVENANCE-BEGIN
% schema: ubt-ai-provenance/v1
% tier: C_working
% ai_assistance: disclosed
% human_review: risk-based
% editorial_responsibility: Ing. David Jaroš
% policy: ../../AI_PROVENANCE.md
% notice: Working material; exhaustive human review is not claimed.
% UBT-AI-PROVENANCE-END

\ifdefined\INCLUDEMODE\else
\documentclass[12pt,a4paper]{article}
\usepackage[a4paper,margin=2.5cm]{geometry}
\usepackage{amsmath,amssymb,amsthm,mathtools}
\usepackage{xcolor}
\usepackage[most]{tcolorbox}
\usepackage{hyperref}

\newtheorem{definition}{Definition}[section]
\newtheorem{lemma}[definition]{Lemma}
\newtheorem{theorem}[definition]{Theorem}
\newtheorem{corollary}[definition]{Corollary}
\newtheorem{remark}[definition]{Remark}

\newtcolorbox{statusbox}[1][]{
  colback=yellow!8!white,
  colframe=orange!70!black,
  arc=3pt,
  boxrule=1pt,
  title=Status,
  #1
}
\newtcolorbox{resultbox}[1][]{
  colback=green!8!white,
  colframe=green!50!black,
  arc=3pt,
  boxrule=1pt,
  title=Result,
  #1
}
\newcommand{\Ltwo}{\textbf{[L2 cand.]}}

\title{Candidate UV Origin of the Electroweak Symmetry-Breaking Sign\\
       \large Compact $\psi$-Sector Casimir Closure for the UBT Alpha Track}
\author{Ing.\ David Jaro\v{s}}
\date{2026-06-13}

\begin{document}
\maketitle
\fi

\begin{abstract}
The previous alpha-closure notes reduced the remaining condition to the sign of
the electroweak quadratic term: the projected low-energy potential must contain
\[
  V_{\rm EW}(\Phi)
  =V_0-\mu_{\rm EW}^2\Phi^\dagger\Phi
  +\lambda_{\rm EW}(\Phi^\dagger\Phi)^2,
  \qquad
  \mu_{\rm EW}^2>0,
  \lambda_{\rm EW}>0.
\]
This note gives the sharpest possible UV closure compatible with the current
repository: if the compact $\psi$-sector induces a negative quadratic Casimir
correction in the minimal electroweak doublet channel, and if the hard bare
quadratic term is absent by modular criticality, then $\mu_{\rm EW}^2>0$ follows.
The proof is algebraic once these two UV hypotheses are accepted.  The note is
therefore a conditional closure theorem, not a claim that the sign has already
been derived unconditionally from the unspecified canonical potential
$V(\Theta)$.
\end{abstract}

\begin{statusbox}
\textbf{Classification.} This is a candidate UV sign theorem.  It closes the
last alpha-proof step under two explicit hypotheses: (H1) modular criticality
forbids a hard positive electroweak mass term at the self-dual point, and (H2)
the compact $\psi$-sector contributes a positive spectral stiffness
$\Gamma_{\psi}>0$ to the symmetry-breaking coefficient.  To promote the alpha
track to fully canonical status, H1 and H2 must be derived from the full
$S[\Theta]$.
\end{statusbox}

\section{Effective quadratic term}

Let $\Phi$ denote the minimal electroweak doublet component of $\Theta$.  The
low-energy radial potential is written as
\[
  V_{\rm EW}(\Phi)
  =V_0+m_{\rm eff}^2\Phi^\dagger\Phi
  +\lambda_{\rm EW}(\Phi^\dagger\Phi)^2+
  O((\Phi^\dagger\Phi)^3).
\]
The symmetry-breaking convention used in the alpha chain is
\[
  m_{\rm eff}^2=-\mu_{\rm EW}^2,
  \qquad
  \mu_{\rm EW}^2>0.
\]
Thus electroweak symmetry breaking is equivalent to
\[
  m_{\rm eff}^2<0.
\]

\section{Compact $\psi$-sector correction}

At the self-dual compact $\psi$ point, write the renormalised quadratic
coefficient in the minimal electroweak doublet channel as
\[
  m_{\rm eff}^2
  =m_{0,\rm EW}^2-\Gamma_{\psi}M_{\psi}^2,
\]
where
\[
  M_{\psi}=R_{\psi}^{-1}
\]
is the compact-fibre scale and $\Gamma_{\psi}$ is the dimensionless spectral
stiffness induced by the $\psi$ modes after Layer2 readout.

The sign convention is chosen so that $\Gamma_{\psi}>0$ means the compact
$\psi$ sector contributes a negative quadratic term to the effective potential,
as in Coleman-Weinberg/Casimir symmetry-breaking mechanisms.

\begin{definition}[Modular criticality hypothesis]
At the self-dual point, no hard electroweak quadratic scale is inserted by hand:
\[
  m_{0,\rm EW}^2=0.
\]
Any non-zero electroweak quadratic term must arise from compact-sector dynamics.
\end{definition}

\begin{definition}[$\psi$-spectral stiffness hypothesis]
The compact $\psi$ sector has positive electroweak-channel stiffness:
\[
  \Gamma_{\psi}>0.
\]
Equivalently, the compact-sector quadratic correction is negative:
\[
  \delta m_{\psi}^2=-\Gamma_{\psi}M_{\psi}^2<0.
\]
\end{definition}

\section{UV sign theorem}

\begin{theorem}[Compact-$\psi$ UV sign closure \Ltwo]
Assume modular criticality,
\[
  m_{0,\rm EW}^2=0,
\]
and positive compact-sector stiffness,
\[
  \Gamma_{\psi}>0.
\]
Then the projected electroweak quadratic coefficient is negative:
\[
  m_{\rm eff}^2<0.
\]
Equivalently, in the convention
\[
  V_{\rm EW}=V_0-\mu_{\rm EW}^2\Phi^\dagger\Phi
  +\lambda_{\rm EW}(\Phi^\dagger\Phi)^2,
\]
one has
\[
  \mu_{\rm EW}^2=\Gamma_{\psi}M_{\psi}^2>0.
\]
\end{theorem}

\begin{proof}
By the effective quadratic decomposition,
\[
  m_{\rm eff}^2
  =m_{0,\rm EW}^2-\Gamma_{\psi}M_{\psi}^2.
\]
Modular criticality gives $m_{0,\rm EW}^2=0$, hence
\[
  m_{\rm eff}^2=-\Gamma_{\psi}M_{\psi}^2.
\]
Since $M_{\psi}^2>0$ and $\Gamma_{\psi}>0$, it follows that
\[
  m_{\rm eff}^2<0.
\]
Therefore
\[
  \mu_{\rm EW}^2=-m_{\rm eff}^2=\Gamma_{\psi}M_{\psi}^2>0.
\]
\end{proof}

\section{Connection to the alpha proof}

Combining this candidate UV sign theorem with the previous low-energy minimum
and vacuum-orbit theorems gives
\[
  \Gamma_{\psi}>0,
  \quad
  m_{0,\rm EW}^2=0
  \Longrightarrow
  \mu_{\rm EW}^2>0
  \Longrightarrow
  \Phi_0\ne0
  \Longrightarrow
  \Theta_0\sim(0,v/\sqrt2),\;Y=1
\]
and hence
\[
  Q_{\rm EM}=T_3+Y/2,
  \qquad
  n=\alpha^{-1}.
\]

\begin{resultbox}
\textbf{Updated alpha status.} With this file, the alpha derivation is complete
as a conditional chain down to two explicit UV hypotheses: modular criticality
and positive $\psi$-sector electroweak stiffness.  The remaining task is no
longer an unspecified ``Higgs gap''; it is the precise derivation of
$\Gamma_{\psi}>0$ and $m_{0,\rm EW}^2=0$ from canonical $S[\Theta]$.
\end{resultbox}

\ifdefined\INCLUDEMODE\else
\end{document}
\fi
